\documentclass[french]{article}

\usepackage{babel}
\usepackage{psfig}
\usepackage{latexsym}
\usepackage{amssymb}
\usepackage{comment}
\usepackage{moreverb}
\usepackage{a4wide}
\usepackage{times}
\usepackage[T1]{fontenc}
\usepackage[french]{minitoc}
\usepackage{latexsym,color}

\newenvironment{solution}{\nopagebreak\par\vspace*{2mm}\underline{Solution~:}\em\par}{\par\vspace*{2mm}}


\newtheorem{defin}{D\'efinition}[section]
\newtheorem{exampl}{Exemple}[section]
\def\picwidth{16cm}

%\excludecomment{solution}

\begin{document}

\vspace*{3cm}

\begin{center}
\Large\bf Contr�le INSTA - Mardi 20 f�vrier 2001
\end{center}

%%%%%%%%%%%%%%%%%%%%%%%%%%%%%%%%%%%%%%%%%%%%%%%%%%%%%%%%%%%%

\vspace*{2cm}

\subsection{Exercices indexation}

Soit un fichier tel que 
chaque page peut contenir 10 articles. On indexe
ce fichier avec un arbre-B,  et on suppose
qu'une page d'index contient 70 paires (cl�, pointeur).
 Les feuilles de l'arbre 
contiennent donc des pointeurs vers le fichier, et les noeuds internes
des pointeurs vers d'autres noeuds. 

\begin{enumerate}
  \item  Donnez (1) le nombre
   de niveaux de l'arbre pour un fichier de 1~000~000 articles, 
   (2) le nombre de
  pages utilis�es (y compris l'index), et (3) le nombre 
d'acc�s disque pour rechercher un article par sa cl�. On suppose que
   toutes les pages sont pleines.


\begin{solution}
Il y a 100~000 pages de donn�es. Un million
d'articles, donc $\frac{1M}{70} = 14~286$ pages pour le dernier
niveau de l'arbre B+. On divise encore par 70\,:
204 pages au niveau $l-1$, 3 pages au niveau sup�rieur, et 1 page
 pour la racine.

Il faut donc lire 4+1 pages pour acc�der � un article
par sa cl�.
\end{solution}

  \item On effectue maintenant une recherche par intervalle
      ramenant 1~000 articles.  D�crivez
    la recherche et donnez le nombre d'acc�s
      disque dans le pire des cas.
\begin{solution}
Recherche sur le plus petit �l�ment (4 pages), puis lecture de 
$\frac{1000}{70} = 14$ pages
en suivant le niveau des feuilles de l'arbre B+. Pour chaque pointeur 
vers un objet il faut lire une page dans le pire des cas, soit
1~000 pages. Total\,: 1018 pages.
\end{solution}


\end{enumerate}
\end{document}
